\documentclass[10pt]{article}

\usepackage[a4paper,margin=20mm]{geometry}
\usepackage{graphicx}
\usepackage{booktabs}
\usepackage{array}
\usepackage{amsmath}
\usepackage{xcolor}
\usepackage{microtype}
\usepackage{newtxtext,newtxmath}
\usepackage{caption}
\usepackage{hyperref}
\usepackage{float}

\definecolor{reviewerblue}{RGB}{0,72,130}
\definecolor{responseblue}{RGB}{0,86,125}
\hypersetup{colorlinks=true,linkcolor=responseblue,urlcolor=responseblue}
\captionsetup{
  font=small,
  labelfont=bf,
  justification=justified,
  singlelinecheck=false
}
\setlength{\textfloatsep}{6pt}
\setlength{\floatsep}{6pt}
\setlength{\intextsep}{6pt}

\newcommand{\assetroot}{reviewer_response_assets/full_v10}
\newcommand{\mechanismroot}{reviewer_response_assets/mechanism}
\newcommand{\reviewercomment}[1]{%
  \noindent\textbf{\textcolor{reviewerblue}{Reviewer comment:}}%
  \ \textit{#1}\par
}
\newcommand{\authorresponse}{%
  \noindent\textbf{\textcolor{responseblue}{Response:}}\ %
}
\newcommand{\pairedplots}[4]{%
  \begin{minipage}[t]{0.485\textwidth}
    \centering
    \textbf{#1}\par\vspace{1mm}
    \includegraphics[width=\linewidth]{\assetroot/#2}
  \end{minipage}\hfill
  \begin{minipage}[t]{0.485\textwidth}
    \centering
    \textbf{#3}\par\vspace{1mm}
    \includegraphics[width=\linewidth]{\assetroot/#4}
  \end{minipage}\par\vspace{2mm}
}

\begin{document}

\section*{Response to the Reviewer: Failure-Case Analysis}

\reviewercomment{Notably, one of the recommendations is to include
failure-case analysis for unsuccessful interception and delayed cooperative
engagement scenarios because understanding framework limitations remains
essential for practical battlefield deployment reliability.}

\authorresponse
Thank you for this important recommendation. We have added a dedicated
failure-case analysis that identifies the terminal miss mechanism, quantifies
the resulting cooperative-engagement delay, and examines the roles of sensing
and actuation imperfections. The analysis uses one reproducible boundary case
from each engagement scenario. These cases are intentionally non-catastrophic:
all eight attacking vehicles are neutralized through assignment redundancy,
but one assigned interceptor fails to enter the \(3~\mathrm{m}\) hit region.
They therefore expose a practically relevant limitation that can be hidden by
target-level success statistics alone.

\paragraph{Boundary-failure outcomes.}
The two examples were obtained with the trained policy fixed under a
one-sample (\(0.05~\mathrm{s}\)) observation delay, position and velocity
measurement-noise standard deviations of \(3~\mathrm{m}\) and
\(0.3~\mathrm{m/s}\), respectively, and the nominal command dynamics. The
evaluation horizon was \(75~\mathrm{s}\), and the cooperative tolerance was
\(0.5~\mathrm{s}\). The reported runs satisfy \(8/8\) target coverage,
\(19/20\) individual interceptor hits, and \(7/8\) complete assignment groups.

\begin{table}[H]
  \centering
  \caption{Observed boundary failures. The miss margin is
  \(d_{\min}-d_{\mathrm{hit}}\), where
  \(d_{\mathrm{hit}}=3~\mathrm{m}\).}
  \label{tab:failure_outcomes}
  \small
  \setlength{\tabcolsep}{4.3pt}
  \renewcommand{\arraystretch}{1.12}
  \begin{tabular}{ccccccccc}
    \toprule
    Case & Seed & Hits & Coverage & Incomplete group &
    \(t_{\mathrm{CPA}}\) (s) & \(d_{\min}\) (m) &
    Margin (m) & Delay lower bound (s) \\
    \midrule
    1 & 76048 & 19/20 & 8/8 & \(A_6\) & 33.10 & 3.243 & 0.243 & \(>41.90\) \\
    2 & 77008 & 19/20 & 8/8 & \(A_2\) & 26.15 & 3.574 & 0.574 & \(>48.90\) \\
    \bottomrule
  \end{tabular}
\end{table}

\paragraph{Direct failure mechanism.}
The incomplete groups arise from different interceptor--target pairs:
\(D_6\!\rightarrow A_6\) in Case~1 and \(D_2\!\rightarrow A_2\) in Case~2.
Both interceptors reach the terminal neighborhood but pass outside the hit
sphere. In Case~1, the \(D_6\)--\(A_6\) range decreases at
\(77.15~\mathrm{m/s}\) one second before closest approach, reaches
\(3.243~\mathrm{m}\) at \(33.10~\mathrm{s}\), and then begins increasing; the
range-rate sign change is observed one integration sample later. In Case~2,
the corresponding values for \(D_2\)--\(A_2\) are
\(79.22~\mathrm{m/s}\), \(3.574~\mathrm{m}\), and \(26.15~\mathrm{s}\).
Thus, the immediate cause is a closest-approach point outside the terminal hit
region, followed by separation, rather than failure to approach the assigned
target.

The terminal geometry is demanding. At closest approach, the line-of-sight
(LOS) angular rates are \(23.89~\mathrm{rad/s}\) and
\(19.31~\mathrm{rad/s}\), and the velocity-to-LOS errors are
\(97.3^\circ\) and \(77.9^\circ\) for Cases~1 and 2, respectively. However,
neither lateral command component saturates during the final two seconds; the
maximum command-vector magnitudes are \(0.648\,g\) and \(0.789\,g\).
Consequently, these failures are better characterized as terminal
phase/geometry errors induced by imperfect state information and command
dynamics, not as insufficient control authority. Case~2 is additionally
exposed to a target acceleration of up to \(6.80~\mathrm{m/s^2}\) during
the final two seconds, whereas the Case~1 target is nearly non-maneuvering in
the same terminal window. This explains the larger miss margin and heading
error in Case~2.

\begin{figure}[H]
  \centering
  \includegraphics[width=\textwidth]
    {\mechanismroot/failure_mechanism_diagnostics_v10.pdf}
  \caption{Terminal-window diagnosis for the missed
  \(D_6\)--\(A_6\) (Case~1) and \(D_2\)--\(A_2\) (Case~2) engagements. The
  distance minimum remains above the \(3~\mathrm{m}\) hit radius in both
  cases. The subsequent change from negative to positive \(\dot d\) confirms
  passage through closest approach and transition to separation. The LOS-rate
  spike reveals the severe terminal geometry, while \(|n_y|\) and \(|n_z|\)
  remain below their component limits.}
  \label{fig:failure_mechanism}
\end{figure}

\paragraph{Why the cooperative engagement is delayed.}
In Case~1, \(D_6\) and \(D_{14}\) are assigned to \(A_6\). \(D_{14}\)
intercepts \(A_6\) at \(33.10~\mathrm{s}\), exactly when \(D_6\) reaches its
\(3.243~\mathrm{m}\) closest point. In Case~2, the redundant
\(D_2,D_{10},D_{18}\) group is assigned to \(A_2\). \(D_{18}\) and
\(D_{10}\) intercept \(A_2\) at \(26.10~\mathrm{s}\) and
\(26.15~\mathrm{s}\), while \(D_2\) passes \(0.574~\mathrm{m}\) outside the
hit boundary at \(26.15~\mathrm{s}\). All other complete groups satisfy the
\(0.5~\mathrm{s}\) tolerance; their maximum hit-time spreads are
\(0.05~\mathrm{s}\) in Case~1 and \(0.35~\mathrm{s}\) in Case~2.

Because the respective missed member never generates a hit event, strict
all-member group completion is right-censored at the \(75~\mathrm{s}\)
horizon. Measured from the first peer hit, the cooperative-completion delay is
therefore greater than \(41.90~\mathrm{s}\) in Case~1 and
\(48.90~\mathrm{s}\) in Case~2.
This distinction is important: the mission-level threat is removed on time,
but the stricter assigned-group objective fails because one member does not
recover after closest-approach passage.

Case~2 also contains a different type of delay in the otherwise complete
\(A_3\) group. Its three members hit at \(64.55\), \(64.80\), and
\(64.90~\mathrm{s}\). Hence \(A_3\) is intercepted much later in absolute
time than the other complete groups, but its within-group spread is only
\(0.35~\mathrm{s}<0.5~\mathrm{s}\). It is therefore a \emph{delayed but
synchronized} cooperative engagement, not a synchronization failure.

\begin{figure}[H]
  \centering
  \includegraphics[width=\textwidth]
    {\mechanismroot/group_absolute_timing_v10.pdf}
  \caption{Absolute group interception times. The dashed orange segments mark
  right-censored incomplete groups (\(A_6\) in Case~1 and \(A_2\) in
  Case~2). The Case~2 \(A_3\) group completes near \(64.9~\mathrm{s}\), much
  later than the other complete groups, although its first-to-last hit spread
  remains within the synchronization tolerance.}
  \label{fig:absolute_group_timing}
\end{figure}

\paragraph{Attribution of sensing and actuation effects.}
To avoid assigning the miss to delay by assumption, we repeated each
representative evaluation with the same frozen policy, seed, assignment, and
target motion while removing one imperfection at a time. Table
\ref{tab:counterfactual} reports individual hits \(H\), complete cooperative
groups \(G\), and the originally missed pair's minimum distance.

\begin{table}[H]
  \centering
  \caption{Frozen-policy counterfactual diagnosis. All conditions retain
  \(8/8\) target coverage. Values are \(H/20\), \(G/8\), and
  \(d_{\min}\) for the originally missed interceptor (m).}
  \label{tab:counterfactual}
  \small
  \setlength{\tabcolsep}{5.0pt}
  \renewcommand{\arraystretch}{1.13}
  \begin{tabular}{lcccccc}
    \toprule
    & \multicolumn{3}{c}{Case 1} & \multicolumn{3}{c}{Case 2} \\
    \cmidrule(lr){2-4}\cmidrule(lr){5-7}
    Condition & \(H\) & \(G\) & \(d_{\min}\) &
                \(H\) & \(G\) & \(d_{\min}\) \\
    \midrule
    Observed boundary condition     & 19 & 7 & 3.243 & 19 & 7 & 3.574 \\
    Remove observation delay only  & 19 & 7 & 3.602 & 17 & 6 & 4.233 \\
    Remove measurement noise only  & 20 & 8 & 1.023 & 20 & 8 & 1.078 \\
    Ideal observation              & 20 & 8 & 1.714 & 20 & 8 & 1.648 \\
    Remove command lag only        & 19 & 7 & 3.432 & 20 & 7 & 2.782 \\
    Ideal observation and no lag   & 20 & 8 & 0.725 & 20 & 8 & 1.599 \\
    \bottomrule
  \end{tabular}
\end{table}

The counterfactual results show two different sensitivity patterns. Case~1 is
a noise-dominated narrow-margin failure. Removing measurement noise moves
\(D_6\) from \(3.243\) to \(1.023~\mathrm{m}\) and restores all eight groups,
whereas removing delay or command lag alone does not recover the hit. Its
\(0.243~\mathrm{m}\) miss margin is much smaller than the \(3~\mathrm{m}\)
position-noise standard deviation, and the ideal-observation conditions also
restore \(20/20\) hits. The slight increase of \(d_{\min}\) when only delay or
lag is removed further confirms that isolated phase changes need not improve
a nonlinear terminal encounter.

Case~2 exhibits a nonlinear interaction. Removing measurement noise restores
\(20/20\) hits and \(8/8\) complete groups, and removing command lag restores
the \(D_2\) hit; however, removing delay alone while retaining noisy
measurements changes the closed-loop phase and produces \(17/20\) hits and
\(6/8\) complete groups. The latter result rules out the simplistic
interpretation that a \(0.05~\mathrm{s}\) observation delay is the sole
failure cause or that performance varies monotonically with delay. Instead,
the data support a coupled mechanism: delayed/noisy terminal state estimates,
command dynamics, high closing speed, and, in Case~2, continuing target
maneuver jointly shift the closest-approach geometry. Since this is a
nonlinear closed loop, the one-factor removals are diagnostic rather than an
additive causal decomposition.

\paragraph{State, control, and cooperative histories.}
Figs.~\ref{fig:failure_trajectory}--\ref{fig:failure_cooperation} provide the
trajectory, acceleration-command, kinematic, attitude, time-to-go,
range-to-go, and group-synchronization histories for the same two runs.

\begin{figure}[p]
  \centering
  \pairedplots{(a) Case 1: planar trajectory}
    {artmappo_case1_seed76048_trajectory.pdf}
    {(b) Case 2: planar trajectory}
    {artmappo_case2_seed77008_trajectory.pdf}
  \pairedplots{(c) Case 1: 3-D trajectory}
    {artmappo_case1_seed76048_trajectory_3d.pdf}
    {(d) Case 2: 3-D trajectory}
    {artmappo_case2_seed77008_trajectory_3d.pdf}
  \caption{Trajectories for the two boundary failures. The emphasized dashed
  curve identifies the missed assigned interceptor; circles denote hit events,
  and the cross denotes its closest-approach point. Redundant assignment
  preserves \(8/8\) target coverage despite the individual miss.}
  \label{fig:failure_trajectory}
\end{figure}

\begin{figure}[p]
  \centering
  \includegraphics[width=0.52\textwidth]
    {\assetroot/artmappo_duav_legend.pdf}\par\vspace{1mm}
  \pairedplots{(a) Case 1: \(n_x\)}
    {artmappo_case1_seed76048_nx.pdf}
    {(b) Case 2: \(n_x\)}
    {artmappo_case2_seed77008_nx.pdf}
  \pairedplots{(c) Case 1: \(n_y\)}
    {artmappo_case1_seed76048_ny.pdf}
    {(d) Case 2: \(n_y\)}
    {artmappo_case2_seed77008_ny.pdf}
  \pairedplots{(e) Case 1: \(n_z\)}
    {artmappo_case1_seed76048_nz.pdf}
    {(f) Case 2: \(n_z\)}
    {artmappo_case2_seed77008_nz.pdf}
  \caption{Acceleration-command histories. The missed member receives bounded
  commands throughout the terminal encounter; its failure is associated with
  the phase and geometry documented in Fig.~\ref{fig:failure_mechanism}, not
  command divergence.}
  \label{fig:failure_controls}
\end{figure}

\begin{figure}[p]
  \centering
  \includegraphics[width=0.52\textwidth]
    {\assetroot/artmappo_duav_legend.pdf}\par\vspace{1mm}
  \pairedplots{(a) Case 1: speed}
    {artmappo_case1_seed76048_velocity.pdf}
    {(b) Case 2: speed}
    {artmappo_case2_seed77008_velocity.pdf}
  \pairedplots{(c) Case 1: heading}
    {artmappo_case1_seed76048_heading.pdf}
    {(d) Case 2: heading}
    {artmappo_case2_seed77008_heading.pdf}
  \pairedplots{(e) Case 1: yaw}
    {artmappo_case1_seed76048_yaw.pdf}
    {(f) Case 2: yaw}
    {artmappo_case2_seed77008_yaw.pdf}
  \pairedplots{(g) Case 1: pitch}
    {artmappo_case1_seed76048_pitch.pdf}
    {(h) Case 2: pitch}
    {artmappo_case2_seed77008_pitch.pdf}
  \caption{Kinematic and attitude histories. Successful interceptors terminate
  near their hit instants, whereas the missed member remains active after
  closest-approach passage.}
  \label{fig:failure_kinematics}
\end{figure}

\begin{figure}[p]
  \centering
  \includegraphics[width=0.52\textwidth]
    {\assetroot/artmappo_duav_legend.pdf}\par\vspace{1mm}
  \pairedplots{(a) Case 1: \(t_{go}\)}
    {artmappo_case1_seed76048_tgo.pdf}
    {(b) Case 2: \(t_{go}\)}
    {artmappo_case2_seed77008_tgo.pdf}
  \pairedplots{(c) Case 1: \(t_{go}\) error}
    {artmappo_case1_seed76048_tgo_error.pdf}
    {(d) Case 2: \(t_{go}\) error}
    {artmappo_case2_seed77008_tgo_error.pdf}
  \pairedplots{(e) Case 1: range-to-go}
    {artmappo_case1_seed76048_distance.pdf}
    {(f) Case 2: range-to-go}
    {artmappo_case2_seed77008_distance.pdf}
  \pairedplots{(g) Case 1: group synchronization}
    {artmappo_case1_seed76048_time_sync.pdf}
    {(h) Case 2: group synchronization}
    {artmappo_case2_seed77008_time_sync.pdf}
  \caption{Cooperative-engagement diagnostics. Successful members terminate
  at interception, whereas the missed member passes its closest point and its
  range subsequently increases. The hatched \(A_6\) (Case~1) and \(A_2\)
  (Case~2) bars report the right-censored lower bounds on strict
  group-completion delay.}
  \label{fig:failure_cooperation}
\end{figure}

\paragraph{Reliability implication and limitation.}
The results establish a two-level conclusion. At the mission level, redundant
assignment prevents a single near-threshold miss from becoming a target
penetration: all eight attackers are intercepted in both examples. At the
strict group level, the current framework does not explicitly detect
closest-approach passage and does not invoke reassignment or a recovery mode
after the assigned target has been hit by peer interceptors. It therefore
cannot guarantee all-member completion under the coupled terminal
uncertainties identified above.

This limitation motivates three concrete extensions: (i) online miss
detection using predicted closest-approach distance together with
\(\dot d\geq 0\), (ii) event-triggered target reassignment or reserve
interceptor activation, and (iii) uncertainty-aware state estimation and
robust training that explicitly cover sensing/actuation phase variations.
The added analysis thus characterizes not only where the present framework
fails, but also which observable signatures can trigger a practical recovery
mechanism.

\end{document}
